\documentclass[10pt,a4paper]{article}

% -------------------- Packages --------------------
\usepackage[margin=0.88in]{geometry}
\usepackage{amsmath,amssymb,mathtools}
\usepackage{booktabs,longtable,tabularx,array,multirow}
\usepackage{graphicx}
\usepackage{float}
\usepackage{xcolor}
\usepackage{enumitem}
\usepackage{caption}
\usepackage{subcaption}
\usepackage{hyperref}
\usepackage{siunitx}
\usepackage{microtype}
% Academic article layout: centered page numbers, no running headers

\hypersetup{
  colorlinks=true,
  linkcolor=blue!55!black,
  urlcolor=blue!55!black,
  citecolor=blue!55!black
}
\sisetup{detect-all,group-separator={,},group-minimum-digits=4}

\setlength{\parskip}{0.22em}
\setlength{\parindent}{1em}
\renewcommand{\arraystretch}{1.15}

\pagestyle{plain}
\setlength{\emergencystretch}{2em}
\renewcommand{\abstractname}{Abstract}
\renewcommand{\arraystretch}{1.12}
\usepackage{titling}
\pretitle{\begin{center}\Large\bfseries}
\posttitle{\end{center}\vspace{-0.4em}}
\preauthor{\begin{center}\normalsize}
\postauthor{\end{center}\vspace{-0.6em}}
\predate{\begin{center}\small}
\postdate{\end{center}\vspace{-0.4em}}


\newcommand{\R}{\mathbb{R}}
\newcommand{\E}{\mathbb{E}}
\newcommand{\Var}{\operatorname{Var}}
\newcommand{\MAE}{\operatorname{MAE}}
\newcommand{\RMSE}{\operatorname{RMSE}}
\newcommand{\MAPE}{\operatorname{MAPE}}
\newcommand{\ME}{\operatorname{ME}}
\newcommand{\logit}{\operatorname{log}}
\newcommand{\placeholder}[1]{\textit{[Insert/verify: #1]}}

\title{Comparative Forecasting of Daily NIFTY 50 Prices Using ARIMA, GARCH, Random Forest, and XGBoost Models}
\author{Shyaam Ganesh\\[-1pt]
\small Department of Mathematics, Indian Institute of Technology Delhi, India}
\date{}

\begin{document}
\maketitle
\begin{abstract}
This study compares classical time-series and machine-learning approaches for forecasting the daily NIFTY 50 closing index. The analysis begins with exploratory analysis and statistical diagnostics of daily log returns, followed by AIC/BIC-based ARMA order selection. Conditional heteroskedasticity is then modelled using GARCH(1,1) errors under Normal and Student-$t$ innovation distributions. Finally, Random Forest and XGBoost regression models are trained using five lagged daily log returns as predictors. Models are evaluated on a common chronological test set of 500 observations. Forecasted returns are converted into closing-price forecasts using the actual previous-day closing price, corresponding to rolling one-step-ahead evaluation. Among the specifications considered, the zero-drift random-walk benchmark, represented by ARIMA(0,1,0), obtains the lowest price-level MAE, RMSE and MAPE. XGBoost is the strongest of the two machine-learning models, while Random Forest has the largest price-level errors among the reported models.
\end{abstract}
\vspace{-0.4em}
\noindent\textbf{Keywords:} NIFTY 50; financial time series; forecasting; ARIMA; GARCH; Random Forest; XGBoost; model comparison
\vspace{0.8em}


% ==================================================
\section{Introduction}
% ==================================================

Forecasting financial time series is important for investment analysis, quantitative decision-making, and risk management. Daily equity-index returns are challenging to predict because their conditional mean may exhibit serial dependence while their volatility can change over time. Linear time-series models provide an interpretable benchmark for mean dynamics, volatility models allow conditional variance to evolve, and machine-learning methods offer a flexible way to represent nonlinear relationships.

This study compares these approaches using daily NIFTY 50 closing prices. The analysis proceeds from exploratory statistical diagnostics to ARIMA order selection, GARCH extensions, and two tree-based machine-learning algorithms. All methods are evaluated on a common chronological test set, allowing their one-step-ahead price forecasts to be compared under the same evaluation design.

The objective is to compare alternative methods for forecasting the daily NIFTY 50 closing index. The modelling sequence is deliberately progressive:
\begin{enumerate}[leftmargin=2em]
  \item inspect the price and return series and assess their statistical properties;
  \item fit ARMA models to daily log returns, selecting orders using AIC and BIC;
  \item extend the mean models with GARCH errors to account for time-varying conditional volatility;
  \item fit Random Forest and XGBoost models to investigate whether nonlinear relationships in lagged returns improve forecasts;
  \item compare all methods on the same final 500-observation test period.
\end{enumerate}

The progression from one model family to the next is motivated by the properties not fully addressed by the preceding model. In particular, evidence of volatility clustering motivates GARCH extensions, while the possibility of nonlinear dependence motivates the machine-learning models.

\section{Methodology}
\subsection{Forecast target and evaluation design}

The models are estimated using daily log returns. Let $P_t$ denote the NIFTY 50 closing price on day $t$. Define
\begin{equation}
  y_t=\log(P_t), \qquad
  r_t=\Delta y_t=\log(P_t)-\log(P_{t-1}).
  \label{eq:returns}
\end{equation}
The target is the ordinary daily log return $r_t$. The first difference of returns, $\Delta r_t=\Delta^2\log(P_t)$, was also examined during preliminary stationarity analysis, but it was not used as the forecasting target. The ARMA and machine-learning models use $r_t$, not $\Delta r_t$.

The data are divided chronologically. The final 500 observations are reserved for out-of-sample testing. For the machine-learning models, the training sample is further divided chronologically into a fitting subset and a validation subset; the final test set is not used for hyperparameter selection.

All reported price forecasts are one-step-ahead forecasts. At each test date, the previous day's observed closing price is used to convert the predicted return to a predicted price:
\begin{equation}
  \widehat{P}_t=P_{t-1}\exp(\widehat{r}_t).
  \label{eq:price_reconstruction}
\end{equation}
Thus, these are rolling one-step-ahead forecasts, not a recursively accumulated 500-step forecast path from a single origin.

% ==================================================
\subsection{Data Collection and Preprocessing}
% ==================================================

Daily NIFTY 50 data were obtained from Yahoo Finance using the ticker \texttt{\^NSEI} and the R package \texttt{quantmod}. The requested start date was 3 November 1995. However, the retrieved Yahoo Finance closing-price series began on 18 September 2007. Consequently, the empirical analysis uses the observations available from that date through the end of the downloaded series.

The closing-price series was transformed into logarithms, and ordinary log returns were obtained by first differencing the log prices as in Equation~\eqref{eq:returns}. Missing observations were removed before constructing the lagged machine-learning dataset.

\begin{table}[H]
\centering
\caption{Data and split used in the machine-learning dataset}
\begin{tabular}{lr}
\toprule
Item & Number of observations \\
\midrule
Lagged machine-learning dataset & 4,690 \\
Training observations & 4,190 \\
Test observations & 500 \\
\bottomrule
\end{tabular}
\label{tab:data_split}
\end{table}

The ARIMA and GARCH models use the same chronological principle: the last 500 observations are held out for testing, with the preceding observations used for estimation.

% ==================================================
\section{Exploratory Analysis and Diagnostic Tests}
% ==================================================

Before estimating forecasting models, the return series was examined using time-series plots, rolling statistics, ACF and PACF plots, and tests for serial dependence and conditional heteroskedasticity.

\subsection{Visual analysis}

The following plots were inspected:
\begin{itemize}[leftmargin=2em]
  \item daily log returns over time;
  \item rolling mean of returns;
  \item rolling standard deviation of returns;
  \item ACF and PACF of returns;
  \item ACF of squared returns.
\end{itemize}

The rolling mean fluctuated around a relatively small value. The rolling standard deviation varied substantially over time, indicating changes in local volatility. The ACF of squared returns showed persistent positive autocorrelation, which is consistent with volatility clustering. These features suggested that modelling the conditional variance could be useful in addition to modelling the conditional mean.

\subsection{Stationarity tests}

ADF, KPSS and Phillips--Perron (PP) tests were conducted during the preliminary analysis. The tests were applied in the investigation of the log-price, ordinary-log-return and differenced-log-return series. Their null hypotheses differ:
\begin{itemize}[leftmargin=2em]
  \item ADF and PP tests have a unit-root null hypothesis;
  \item the KPSS test has a stationarity null hypothesis.
\end{itemize}

The final forecasting target was retained as the ordinary log return $r_t=\Delta\log(P_t)$. The first-differenced return series was not substituted as the target merely because an additional difference produced more favourable stationarity-test outcomes. For a formal final submission, the numerical ADF, KPSS and PP statistics and the exact input series for each test should be copied from the original R output and reported together. This is important because the test results discussed during model development were not consistently labelled by input series in the working record.

\subsection{Serial correlation in returns}

The Ljung--Box test was applied to raw returns. The reported results were:
\begin{table}[H]
\centering
\caption{Ljung--Box tests on raw daily log returns}
\begin{tabular}{rrr}
\toprule
Lag & Test statistic & $p$-value \\
\midrule
10 & 41.435 & $9.451\times 10^{-6}$ \\
20 & 72.265 & $7.740\times 10^{-8}$ \\
40 & 104.67 & $1.078\times 10^{-7}$ \\
\bottomrule
\end{tabular}
\label{tab:raw_lb}
\end{table}

The null hypothesis of no autocorrelation was rejected at each reported lag. This provided a motivation for fitting ARMA mean equations rather than relying exclusively on a zero-mean return model.

\subsection{Motivation for GARCH modelling}

The ACF of squared returns showed strong dependence, and the rolling standard deviation varied over time. In addition, ARCH-LM testing was used to investigate conditional heteroskedasticity. These findings motivated the extension from ARMA models to ARMA--GARCH models. GARCH modelling is intended primarily to capture conditional variance dynamics; it does not guarantee a material improvement in point forecasts of the conditional mean or closing price.

% ==================================================
\section{ARIMA Model Selection and Forecasting}
% ==================================================

\subsection{Candidate models and selection criteria}

ARMA$(p,q)$ models were estimated on the ordinary log-return series. A grid search considered
\begin{equation}
  p,q\in\{0,1,\ldots,10\},
\end{equation}
and considered specifications with and without a mean term. The candidate models were compared using the Akaike Information Criterion (AIC) and Bayesian Information Criterion (BIC). Lower values are preferred within the same estimation sample and likelihood framework.

The selected specifications were:
\begin{itemize}[leftmargin=2em]
  \item AIC selected ARMA$(5,5)$ with a mean term, equivalent to ARIMA$(5,1,5)$ with drift when expressed in log prices;
  \item BIC selected ARMA$(0,0)$ without a mean term, equivalent to a zero-drift ARIMA$(0,1,0)$ model on log prices.
\end{itemize}

\begin{table}[H]
\centering
\caption{Selected ARIMA models and information criteria}
\begin{tabular}{lrr}
\toprule
Model & AIC & BIC \\
\midrule
ARIMA$(5,1,5)$ & $-24{,}359.43$ & $-24{,}283.33$ \\
ARIMA$(0,1,0)$ & $-24{,}308.52$ & $-24{,}302.18$ \\
\bottomrule
\end{tabular}
\label{tab:arima_ic}
\end{table}

AIC favoured the more complex ARIMA$(5,1,5)$, whereas BIC favoured the parsimonious ARIMA$(0,1,0)$. The two criteria therefore provided complementary model choices rather than a single unanimous selection.

\subsection{Residual checking and motivation for GARCH}

Residual diagnostics were examined after fitting the selected ARMA models. For the AIC-selected ARMA$(5,5)$ model, a Ljung--Box test at lag 20, using 10 fitted degrees of freedom, gave a test statistic of 23.381 and a $p$-value of 0.009426. This indicates that some residual serial dependence remained. The BIC-selected ARMA$(0,0)$ model also showed autocorrelation in its residual ACF.

Separately, the squared-return diagnostics suggested strong volatility dependence. This distinction motivated the next modelling stage: ARMA terms were retained to describe the conditional mean, while GARCH terms were introduced to describe the time-varying conditional variance. The GARCH extension was therefore motivated by the volatility behaviour, not by an assumption that GARCH would necessarily improve the point forecast of the index level.

\subsection{ARIMA out-of-sample forecasts}

Rolling one-step-ahead forecasts were generated over the final 500 observations. Predicted returns were converted to price forecasts using Equation~\eqref{eq:price_reconstruction}. The resulting price-level metrics are reported in Section~\ref{sec:overall_results}.

% ==================================================
\section{ARMA--GARCH Models}
% ==================================================

\subsection{Model formulation}

For an ARMA$(p,q)$ conditional mean and a GARCH$(1,1)$ conditional variance, the specification is
\begin{align}
  r_t &= \mu + \sum_{i=1}^{p}\phi_i r_{t-i}
       + \sum_{j=1}^{q}\theta_j\epsilon_{t-j}+\epsilon_t, \\
  \epsilon_t &= \sigma_t z_t, \\
  \sigma_t^2 &= \omega+\alpha_1\epsilon_{t-1}^2+\beta_1\sigma_{t-1}^2.
\end{align}
Here $z_t$ follows either a standard Normal or a standardised Student-$t$ distribution. The parameter $\alpha_1$ captures the response of current variance to the previous squared innovation, while $\beta_1$ captures persistence from the previous conditional variance.

Four specifications were estimated: ARMA$(5,5)$--GARCH$(1,1)$ with Normal innovations, ARMA$(5,5)$--GARCH$(1,1)$ with Student-$t$ innovations, and the corresponding two ARMA$(0,0)$--GARCH$(1,1)$ specifications. The ARMA$(0,0)$ models were specified without a mean term to match the zero-drift random-walk benchmark.

\subsection{Information-criterion comparison}

\begin{table}[H]
\centering
\caption{Information criteria for the four ARMA--GARCH specifications}
\small
\begin{tabular}{lrrr}
\toprule
Model & Log-likelihood & AIC & BIC \\
\midrule
ARMA$(5,5)$--GARCH Normal & 13153.21 & -6.2627 & -6.2416 \\
ARMA$(5,5)$--GARCH Student-$t$ & 13249.03 & -6.3079 & -6.2853 \\
ARMA$(0,0)$--GARCH Normal & 13124.05 & -6.2541 & -6.2495 \\
ARMA$(0,0)$--GARCH Student-$t$ & 13215.71 & -6.2973 & -6.2912 \\
\bottomrule
\end{tabular}
\label{tab:garch_ic}
\end{table}

Within both conditional-mean specifications, Student-$t$ innovations produced lower AIC and BIC than Normal innovations. Thus, the Student-$t$ distribution was preferred by the information criteria. This choice is consistent with the possibility of heavier tails in financial returns, although the goodness-of-fit diagnostics discussed below still rejected the fitted innovation distributions.

\subsection{Estimated variance parameters}

\begin{table}[H]
\centering
\caption{Estimated GARCH variance parameters}
\begin{tabular}{lrrr}
\toprule
Model & $\alpha_1$ & $\beta_1$ & $\alpha_1+\beta_1$ \\
\midrule
ARMA$(5,5)$ Normal & 0.100961 & 0.891103 & 0.992064 \\
ARMA$(5,5)$ Student-$t$ & 0.079598 & 0.910050 & 0.989648 \\
ARMA$(0,0)$ Normal & 0.094073 & 0.898206 & 0.992279 \\
ARMA$(0,0)$ Student-$t$ & 0.076071 & 0.913095 & 0.989166 \\
\bottomrule
\end{tabular}
\label{tab:garch_parameters}
\end{table}

All reported persistence measures, $\alpha_1+\beta_1$, are close to one and below one. This indicates highly persistent estimated conditional volatility. It should not be interpreted as evidence that the parameters are constant over the full sample, because the stability tests below reject parameter stability.

\subsection{Residual and specification diagnostics}

The diagnostics used to assess the fitted ARMA--GARCH models were:
\begin{enumerate}[leftmargin=2em]
  \item Ljung--Box tests on standardised residuals, to check for remaining serial correlation in the conditional mean;
  \item Ljung--Box tests on squared standardised residuals, to check for remaining dependence in the conditional variance;
  \item ARCH-LM tests, to check for remaining ARCH effects;
  \item Nyblom stability tests, to assess parameter constancy;
  \item sign-bias tests, to investigate whether positive and negative shocks have different volatility effects;
  \item adjusted Pearson goodness-of-fit tests, to assess the assumed innovation distribution.
\end{enumerate}

\begin{table}[H]
\centering
\caption{Summary of ARMA--GARCH diagnostic results}
\small
\begin{tabularx}{\textwidth}{p{0.28\textwidth}XXXX}
\toprule
Diagnostic & ARMA$(5,5)$ Normal & ARMA$(5,5)$ Student-$t$ & ARMA$(0,0)$ Normal & ARMA$(0,0)$ Student-$t$ \\
\midrule
Standardised residual Ljung--Box & Pass & Pass & Fail & Fail \\
Squared residual Ljung--Box & Pass & Pass & Pass & Pass \\
ARCH-LM & Pass & Pass & Pass & Pass \\
Nyblom stability & Fail & Fail & Fail & Fail \\
Sign bias / joint effect & Fail & Fail & Fail & Fail \\
Adjusted Pearson goodness-of-fit & Fail & Fail & Fail & Fail \\
\bottomrule
\end{tabularx}
\label{tab:garch_diagnostics}
\end{table}

For the ARMA$(5,5)$ models, the standardised-residual Ljung--Box tests did not indicate remaining serial correlation at the reported lags. Both ARMA$(0,0)$ models failed these tests, indicating that the zero-mean conditional-mean equation left serial dependence unexplained. All four models passed the reported squared-residual Ljung--Box and ARCH-LM tests, suggesting that the GARCH variance equation captured much of the remaining volatility dependence.

However, all four models failed the Nyblom stability and sign-bias diagnostics, and the adjusted Pearson tests rejected the assumed Normal or Student-$t$ innovation distributions. Therefore, no ARMA--GARCH specification passed every diagnostic. The Student-$t$ ARMA$(5,5)$ model was the information-criterion favourite, but the diagnostic evidence calls for caution in interpreting it as a fully adequate model.

\subsection{Why proceed to machine learning?}

The ARMA--GARCH analysis provides a structured parametric description of conditional mean and variance dynamics. Nevertheless, the remaining sign bias and parameter-instability evidence suggest that the fitted specifications may not capture all aspects of the return-generating process. In addition, the ARMA$(0,0)$ specification leaves serial correlation in standardised residuals. Random Forest and XGBoost were therefore introduced as alternative flexible regression methods to investigate whether nonlinear relationships among lagged returns could improve out-of-sample point forecasts. This is an empirical comparison, not an assumption that machine-learning methods must perform better.

% ==================================================
\section{Machine-Learning Forecasting}
% ==================================================

\subsection{Supervised-learning formulation}

The machine-learning target is the ordinary daily log return $r_t$. Five lagged returns are used as predictors:
\begin{equation}
  X_t=(r_{t-1},r_{t-2},r_{t-3},r_{t-4},r_{t-5}),
  \qquad Y_t=r_t.
  \label{eq:ml_features}
\end{equation}
The supervised-learning problem is to estimate
\begin{equation}
  \widehat r_t=f(X_t).
\end{equation}
The predicted return is converted into a closing-price forecast using Equation~\eqref{eq:price_reconstruction}. Each test observation uses the five actual preceding returns available at that forecast date, corresponding to a rolling one-step-ahead setting.

The chronological training sample was split again into a fitting subset (80\%) and a validation subset (20\%). Hyperparameters were selected using validation RMSE. The final 500-observation test set was kept separate from tuning.

\subsection{Random Forest}

Random Forest regression was fitted using the \texttt{randomForest} package. A grid search considered
\begin{itemize}[leftmargin=2em]
  \item \texttt{mtry} $\in\{1,2,3,4,5\}$;
  \item \texttt{nodesize} $\in\{1,5,10\}$;
  \item \texttt{ntree} fixed at 500.
\end{itemize}
This produced 15 configurations. The configuration with the lowest validation RMSE was selected.

\begin{table}[H]
\centering
\caption{Selected Random Forest configuration}
\begin{tabular}{lr}
\toprule
Hyperparameter & Selected value \\
\midrule
\texttt{ntree} & 500 \\
\texttt{mtry} & 1 \\
\texttt{nodesize} & 10 \\
Validation RMSE & 0.009063275 \\
\bottomrule
\end{tabular}
\label{tab:rf_tuning}
\end{table}

The final model was refitted using the entire training sample. Its out-of-bag mean squared residual was 0.0001845519 and its reported variance explained was $-3.72\%$. A negative out-of-bag variance-explained value indicates that the out-of-bag predictions were worse than the corresponding mean-response benchmark under that measure.

The percentage increase in MSE importance values were: Lag1 = 0.6334, Lag2 = 3.1419, Lag3 = 3.2030, Lag4 = 3.2032 and Lag5 = 5.1366. Lag5 ranked highest by this measure. The relatively small values should be interpreted cautiously; they do not establish strong predictive relationships.

\subsection{XGBoost}

XGBoost regression was fitted using the \texttt{xgboost} package. The search considered tree depth, learning rate and minimum child weight. The search was expanded to investigate deeper trees rather than assuming that a shallow tree was optimal. The validation RMSE was calculated directly for each candidate model.

The expanded grid included
\begin{align*}
  \texttt{max\_depth} &\in\{1,2,3,4,5,7,10\},\\
  \texttt{eta} &\in\{0.03,0.05\},\\
  \texttt{min\_child\_weight} &\in\{1,5,10\},\\
  \texttt{nrounds} &\in\{50,100,200\}.
\end{align*}
The expanded search selected depth 2, learning rate 0.03 and minimum child weight 5. A follow-up search examined boosting rounds from 1 through 10. One round yielded the lowest validation RMSE, but the five-round configuration was selected for the final model because it was the user's preferred second-best configuration and its validation performance was very close to the minimum.

\begin{table}[H]
\centering
\caption{Final selected XGBoost configuration}
\begin{tabular}{lr}
\toprule
Hyperparameter & Selected value \\
\midrule
\texttt{max\_depth} & 2 \\
\texttt{eta} & 0.03 \\
\texttt{min\_child\_weight} & 5 \\
\texttt{nrounds} & 5 \\
Validation RMSE & 0.008728660 \\
\bottomrule
\end{tabular}
\label{tab:xgb_tuning}
\end{table}

The final XGBoost model was refitted on the entire training dataset and evaluated on the untouched test set. The choice of five rounds is a deliberate selection among near-equivalent validation results, rather than the absolute minimum-validation-RMSE choice.

\subsection{Why compare Random Forest and XGBoost?}

Both algorithms can represent nonlinear relationships and interactions without requiring a pre-specified linear ARMA structure. Random Forest averages predictions over bootstrap-based decision trees, while XGBoost builds trees sequentially to reduce the loss. Using the same five lagged-return features and the same chronological split provides a controlled comparison of these two tree-based approaches with the classical models.

% ==================================================
\section{Out-of-Sample Forecast Evaluation}
\label{sec:overall_results}
% ==================================================

\subsection{Evaluation metrics}

Let $P_t$ be the actual closing price and $\widehat P_t$ its forecast. The error is defined as $e_t=P_t-\widehat P_t$. The reported price-level metrics are
\begin{align}
  \ME &= \frac{1}{n}\sum_{t=1}^{n}(P_t-\widehat P_t),\\
  \MAE &= \frac{1}{n}\sum_{t=1}^{n}|P_t-\widehat P_t|,\\
  \RMSE &= \sqrt{\frac{1}{n}\sum_{t=1}^{n}(P_t-\widehat P_t)^2},\\
  \MAPE &= \frac{100}{n}\sum_{t=1}^{n}
  \left|\frac{P_t-\widehat P_t}{P_t}\right|.
\end{align}
Lower MAE, RMSE and MAPE indicate smaller forecast errors. Under this sign convention, negative ME indicates average overprediction, while positive ME indicates average underprediction.

\subsection{Consolidated out-of-sample results}

\begin{table}[H]
\centering
\caption{Price-level forecasting performance on the 500-observation test set}
\small
\begin{tabular}{lrrrr}
\toprule
Model & ME & MAE & RMSE & MAPE (\%) \\
\midrule
ARIMA$(5,1,5)$ & -8.2379 & 154.5298 & 205.2892 & 0.6412 \\
ARIMA$(0,1,0)$ & 1.2911 & 151.6817 & 202.3063 & 0.6290 \\
ARMA$(5,5)$--GARCH Normal & -16.1375 & 154.1686 & 203.7747 & 0.6397 \\
ARMA$(5,5)$--GARCH Student-$t$ & -17.3660 & 153.9794 & 203.7966 & 0.6388 \\
ARMA$(0,0)$--GARCH Normal & 1.2911 & 151.6817 & 202.3063 & 0.6290 \\
ARMA$(0,0)$--GARCH Student-$t$ & 1.2911 & 151.6817 & 202.3063 & 0.6290 \\
Tuned Random Forest & -5.4442 & 156.5184 & 205.2712 & 0.6488 \\
Tuned XGBoost & -7.9509 & 152.2220 & 202.5839 & 0.6314 \\
\bottomrule
\end{tabular}
\label{tab:final_results}
\end{table}

The ARIMA$(0,1,0)$ zero-drift benchmark has the lowest MAE, RMSE and MAPE among the reported models. XGBoost is the stronger of the two machine-learning models and is close to the ARIMA$(0,1,0)$ benchmark. Random Forest has the largest MAE, RMSE and MAPE in this comparison.

The ARMA$(0,0)$--GARCH models produce the same point-price forecasts as the zero-drift ARIMA$(0,1,0)$ benchmark because their conditional mean was constrained to zero. The Normal and Student-$t$ variance specifications change the conditional volatility forecasts, but not the zero conditional-mean forecast. Therefore, their price-level point forecast metrics are identical.

\subsection{Interpretation and model-selection conclusion}

The results show that greater model complexity did not automatically improve point-forecast accuracy. The simple zero-drift random-walk benchmark performed best on the selected test period. Among the machine-learning methods, XGBoost performed better than Random Forest on MAE, RMSE and MAPE. The ARMA$(5,5)$--GARCH Student-$t$ model had the best GARCH information criteria, but its diagnostic tests revealed parameter instability, sign bias and lack of fit under the assumed innovation distribution.

The best model therefore depends on the forecasting objective:
\begin{itemize}[leftmargin=2em]
  \item for closing-price point forecasts on this test period, ARIMA$(0,1,0)$ (equivalently, the zero-mean return model) performed best;
  \item for the machine-learning comparison, tuned XGBoost performed best;
  \item for modelling conditional volatility, the GARCH specifications captured much of the volatility clustering, although stability and asymmetry diagnostics indicate that further variance-model work may be warranted.
\end{itemize}

% ==================================================
\section{Actual and Predicted Price Plots}
% ==================================================

The following figure environments include the plots generated during the R analysis. Place the image files in the same directory as this \LaTeX{} source. If a file is not present, the document will compile and display a placeholder instead.

\begin{figure}[H]
\centering
\includegraphics[width=0.85\textwidth]{NIFTY50_ARIMA_Actual_vs_Predicted.png}
\caption{Actual versus predicted NIFTY 50 closing prices using ARIMA models.}
\label{fig:arima}
\end{figure}

\begin{figure}[H]
\centering
\includegraphics[width=0.85\textwidth]{NIFTY50_GARCH_Actual_vs_Predicted.png}
\caption{Actual versus predicted NIFTY 50 closing prices using ARMA--GARCH models.}
\label{fig:garch}
\end{figure}

\begin{figure}[H]
\centering
\includegraphics[width=0.85\textwidth]{NIFTY50_Tuned_RF_Actual_vs_Predicted.png}
\caption{Actual versus predicted NIFTY 50 closing prices using tuned Random Forest.}
\label{fig:randomforest}
\end{figure}

\begin{figure}[H]
\centering
\includegraphics[width=0.85\textwidth]{NIFTY50_Tuned_XGBoost_Actual_vs_Predicted.png}
\caption{Actual versus predicted NIFTY 50 closing prices using tuned XGBoost.}
\label{fig:xgboost}
\end{figure}

% ==================================================
\section{Limitations and Further Work}
% ==================================================

The conclusions are specific to the downloaded Yahoo Finance series and the selected 500-observation test period. Several additional checks would strengthen the analysis:
\begin{enumerate}[leftmargin=2em]
  \item Re-run and document the ADF, KPSS and PP tests with explicit labels identifying whether each test uses log prices, ordinary log returns or first-differenced returns. Report the test statistics, critical values and $p$-values together.
  \item Consider asymmetric volatility specifications such as GJR-GARCH or EGARCH in response to the sign-bias results.
  \item Evaluate volatility forecasts separately using a realised-volatility proxy. Price-level MAE and RMSE do not measure the accuracy of conditional variance forecasts.
  \item Examine robustness across alternative chronological test windows, rather than relying on a single 500-observation period.
  \item Investigate additional predictive features for the machine-learning models, such as lagged squared returns and rolling volatility, while ensuring that all features are constructed using information available at the forecast date.
\end{enumerate}

A further methodological caution is that AIC/BIC values from different likelihood formulations should not be compared indiscriminately. The ARMA information criteria and the GARCH information criteria in this report are reported within their respective model families; the out-of-sample price-level metrics provide the common comparison across all approaches.

% ==================================================
\section{Conclusion}
% ==================================================

This experiment compared ARIMA, ARMA--GARCH, Random Forest and XGBoost methods for daily NIFTY 50 closing-price forecasting. Preliminary diagnostics identified serial dependence and volatility clustering in returns, motivating ARMA modelling and subsequent GARCH extensions. AIC and BIC selected different ARIMA specifications, while Student-$t$ innovations improved the information criteria for both GARCH mean-model families. Diagnostic tests nevertheless indicated instability and remaining specification issues in the GARCH models.

The machine-learning models used five lagged daily returns. Random Forest was tuned over the number of predictors considered at each split and terminal-node size. XGBoost was tuned over tree depth, learning rate, child weight and boosting rounds. On the common 500-observation test set, the zero-drift ARIMA$(0,1,0)$ benchmark recorded the lowest closing-price MAE, RMSE and MAPE. XGBoost was the best-performing machine-learning method and outperformed Random Forest. The results illustrate that, for this series and test period, a simple random-walk benchmark remained difficult to beat in point forecasting, even when more flexible models were introduced.

\end{document}